% Extracted from the compiled manuscript.
% Source label: tab:app_forecast_targets

\begin{table}[!htbp]
\centering
\caption{Realised-variation quantities and forecast-horizon specification.}
\label{tab:app_forecast_targets}
\small
\begin{tabular}{@{}p{0.24\linewidth} p{0.47\linewidth} p{0.22\linewidth}@{}}
\toprule
Quantity & Definition & Role\\
\midrule
Realised variance and volatility
& \(\mathrm{RV}_t=\sum_i r_i^2\) and
\(\mathrm{RVol}_t=\sqrt{\mathrm{RV}_t}\), using within-session log-mid-price
returns with the first session return set to zero
& Forecast target and persistence benchmark\\

Bipower variation
& \(\mathrm{BV}_t=(\pi/2)\sum_i |r_i||r_{i-1}|\)
& Jump-robust continuous-variation estimate\\

Jump component
& \(J_t=\max\{\mathrm{RV}_t-\mathrm{BV}_t,0\}\)
& Non-negative discontinuity measure\\

Forward horizons
& Next complete native window; 30, 45, and 60 seconds; 2, 5, 10, 15, and
30 minutes; and 1 hour
& Multi-horizon forecasting\\

Constraints
& Targets remain within the same trading session, and no forecast horizon is
finer than the native covariance duration
& Temporal-alignment and leakage control\\
\bottomrule
\end{tabular}
\begin{minipage}{0.94\linewidth}\footnotesize
These quantities are constructed separately from the 15 event-level inputs to
\(\Sigma_t\). Realised variance enters the jump decomposition, while realised
volatility supplies the volatility-scale forecasting outcome.
\end{minipage}
\end{table}
